\documentclass{article}
\usepackage[utf8]{inputenc}
\usepackage[T1]{fontenc}
\usepackage[margin=2cm]{geometry}
\usepackage{hyperref}
\sloppy
\begin{document}
\section*{Historical source}
Encyclopaedia Britannica scan, source row 2296. Parse 3.5 extracted the text. Super 3.5 VL corrected it from an enlarged crop of each block. It can still contain recognition errors.
Page headers, footers, and page numbers are removed.

\subsection*{Untitled continuation (E0213)}
\textit{Continues an entry from an earlier page.}

Po.

Ankudinoff's vessel was wrecked, and the crew distributed among the other two. Soon after this the two vessels lost sight of each other, and never joined again. Dishneff was driven about by tempestuous winds till October, when he was shipwrecked considerably to the fourth of the Anadyr. Having at last reached that river, he formed a scheme of returning by the same way that he had come, but never made the attempt. As for Alexeeff, after being also shipwrecked, he had died of the scurvy, together with Ankudinoff's part of the crew were killed by the favages, and a few escaped to Kamschatka, where they settled.

See
Cook's
Difco-
ries, 95.
100.

Kamfchatka, where they settled."
From Captain Cook's voyage towards the north-
eastern parts of Afia, it appears that it is pofible to
double the promontory of Tichutki without any great
difficulty: and it now appears, that the continents of

- icon.

Insurmountable obstacles in the north-east passage.

Of the
south-west
passage.

See
Cook's
Difection-
ries, n° 21
13
Parry's
argu-
ments in
savour of
a possibility
of reaching
the pole.

Of the
south-west
passage.

See
Cook's
Difection-
ries, no. 21
3
Parry's
arguments in
savour of
a possibility
of reaching
the pole.

have reached the billions of 88.--a. From the telli-
veyor of Captain Choyer, who gave anders to certain
queries drawn up by Mr Dalymple concerning the po-
utes first appears to be lead when in the latitude of 83". One Mr Watt informed our author, then when he was 7 years of age, as that time making his first voyage on Captain M McCallon, a le hold an fullina n-
vigor, who commanded a Scotch whale-fishing dip, as during the time that the whales are fupped to copia the so fishing one, the carried on, the captain evolved to employ that interIn in attempting to resta the north pole. He ascondingly revealed because the lead ob-
duction on 88; when the fea was not only open to the northwest, but they had now no ice on the last three degrees, but which he fill advanced, the ante complained that, the compils was not Ready, and the Captain was often dealt with collection to give over his attempts – TD Compell, the combination of Harris’s voyages, informed Mr Harrington, that Dr Jitalie, an-
vivord Mullan, being in his work on board. In March
July of, or which at that time was ufully feat to fu-
penined the Greenland fishery, the captain determined, like the toocimshn above mentioned, to make an attempt to reach the pole during the intervent leaves the drift too forest fisheries. He expressed, according to the belt di Dr Campbell's recollection, as far as 88'3 when the weelwe was warm, the for free from ice, and falling like the bay of Bicky. Zullies now prefilled the expense to proceed:) but he answered, that he had already gone to for, and finally the himself it Holland has regarding his history upon which account he would differ on journal to the kept, but returned is soon as probable to Philology—— as the year (66-20), Mr Dobberg, than fercerty of the Royal Society, was ordered to regite a paper, entitled " Several linguists concerning Greenland, minnows by Mr Gary, who had visited their posts." The 19th of their queries is the following: If we were hath any one been known to approach the plot? —The name is, “I once met upon the coult of Greatlanda Holderar that force me had been half a degree from the days, though one are his journal, which was sila attested by his mares, where they had ten an ice on day, but ali winter” – Go. In Captain Woods’s ac- count of a voyage its quell of the north-ent pulgage, we have the following account of a Dutch high which ended the mistake the 69/o". Captain Condekly, who had made above 30 voyages to Greenland, did relate to his majdly, the being a Greatland form 30 years before, he was at company with two Filhamalians to the ealiward of Edge’s lineal and that the whales are ap- peemings on the shore, the Hindians were determined to be further encountered I and in a fortnight’s time return- ed, and gave it out that they had failed into the life- undys of, and that they did not meet him any face, but a free and open teas, and that there aren a very bullazing groups fail. Next, the job Bicky, Mr Gould, we bring not finished with the bare relation, they pro-duced him four journals of the two ships, which fulfilled them, none, that they agreed them four instances."--7. In the Philadelphia Traffletions for 167; we knew the following pulgage: «For it is well known to all that fish northward, instead of the north-ern coasts are frozen up for many fingers, though in the open fea it is not by no more under the pole right, unless by accident!" In which pulgage, we have the following requestout of 2 Dutch high which could see the ultimate of 69/o". Captain Thomas Robertson offand our author, that he had been in latitude 834, that the fea real open, and he was certain that he could

have reached the Latitude of 83°.—2. From the testimony of Captain Cheyne, who gave answers to certain queries drawn up by Mr Dalrymple concerning the polar seas, it appears that he had been in the latitude of 82°.—3. One Mr Watt informed our author, that when he was 17 years of age, at that time making his first voyage with Captain M'Callam, a bold and skilful navigator, who commanded a Scotch whale-fishing ship, as during the time that the whales are fupplied to copulate no fishing can be carried on, the captain resolved to employ that interval in attempting to reach the north pole. He accordingly proceeded without the least obstruction to 83°, when the sea was not only open to the northward, but they had seen no ice for the last three degrees; but while he still advanced, the mate complained that the compass was not steady, and the captain was obliged with reluctance to give over his attempt.—4. Dr Campbell, the continuator of Harris's voyages, informed Mr Barrington, that Dr Dalie, a native of Holland, being in his youth on board a Dutch ship of war which at that time was usually fitted to superintend the Greenland fishery, the captain determined, like the Scotchman above mentioned, to make an attempt to reach the pole during the interval between the first and second fisheries. He penetrated, according to the belt of Dr Campbell's recollection, as far as 88°; when the weather was warm, the sea free from ice, and rolling like the bay of Biscay. Dalle now pressed the captain to proceed; but he answered, that he had already gone too far, and should be blamed in Holland for neglecting his station; upon which account he would suffer no journal to be kept, but returned as soon as possible to Spitbergen.—5. In the year 1662-3, Mr Oldenburg, then secretary of the Royal Society, was ordered to register a paper, entitled "Several Inquiries concerning Greenland, answered by Mr Gray, who had visited those parts." The 19th of these queries is the following: How near hath any one been known to approach the pole?—The answer is, "I once met upon the coast of Greenland a Hollander that swore he had been half a degree from the pole, showing me his journal, which was also attested by his mate; where they had seen no ice or land, but all water."—6. In Captain Wood's account of a voyage in quest of the north-east passage, we have the following account of a Dutch ship which reached the latitude of 89°. "Captain Goulden, who had made above 30 voyages to Greenland, did relate to his majesty, that being at Greenland some 20 years before, he was in company with two Hollanders to the eastward of Edge's island; and that the whales not appearing on the shore, the Hollanders were determined to go farther northward; and in a fortnight's time returned, and gave it out that they had sailed into the latitude 89°, and that they did not meet with any ice, but a free and open sea, and that there run a very hollow sea like that of the Bay of Biscay. Mr Goulden being not satisfied with the bare relation, they produced him four journals out of the two ships, which testified the same, and that they all agreed within four minutes."—7. In the Philosophical Transactions for 1675 we have the following passage: "For it is well known to all that fail northward, that most of the northern coasts are frozen up for many leagues, though in the open sea it is not so, no nor under the pole itself, unless by accident." In which passage the having reached the


\section*{Source record}
\begin{description}
\item[Scan ID] \texttt{britannica:309f765b1ecc83e0a7f7762e44b4a927a7555443:train:2296}
\item[Image file] \texttt{page-2296-487afb796c01.jpg}
\item[Image SHA-256] \texttt{487afb796c010b5708a60a798cd2b38cff43eea65fd8949d3ae509d3a6f63a13}
\item[Dataset] NationalLibraryOfScotland/encyclopaedia\_britannica\_illustrated at revision \texttt{309f765b1ecc83e0a7f7762e44b4a927a7555443} (CC0-1.0)
\item[Extraction] model \texttt{nvidia/nemotron-3.5-parse}, raw reply \texttt{extract/raw/page-2296-487afb796c01.json}
\item[Raw reply SHA-256] \texttt{f451c7458b1f0fbb2b5849271c4424468d446fa304396f61e00a40f078a6d5d4}
\item[Correction] model \texttt{nvidia/Nemotron-3.5-Super-VL}, block crops, raw replies \texttt{correct/raw/crop\_nothink/page-2296-487afb796c01-b*.json}; parse\_fallback: 2, super\_vl: 9
\item[Running header] **POL [ 93 ]**
\end{description}
\end{document}
